\chapter{Az online vizsgáztatás szakirodalmi háttere}
\label{ch:szakirodalom}
% Terjedelem: kb. 10–12 oldal (~18 000 karakter)
% 13. pont: a dolgozat 50–60%-a szakirodalomra épüljön -> ez és a 3. fejezet adja a gerincét.
% Minden állításnál hivatkozás: \cite{kulcs}, szó szerinti idézetnél: \cite[12.~o.]{kulcs}

\section{Számítógéppel támogatott értékelés}
% E-learning, LMS-ek, a számítógépes értékelés (CAA) fogalma, előnyei/hátrányai.
% Formatív vs. szummatív értékelés. Források: \cite{bull2004blueprint}
\todo{Szöveg}

\section{Feladattípusok és automatikus pontozásuk}
% Zárt (választós, igaz–hamis), félig nyílt (rövid válasz), nyílt (esszé) feladatok.
% Több jó válaszos feladatok pontozási módszerei (részpontozás, levonás a rossz jelölésért).
\todo{Szöveg}

\section{Programozási feladatok automatikus értékelése}
% Teszteset-alapú értékelés, nyilvános és rejtett tesztek, ismert rendszerek (pl. CodeRunner,
% DOMjudge, Web-CAT). Források: \cite{douce2005automatic,ihantola2010review,paiva2022automated}
\todo{Szöveg}

\section{Online felügyelet (proctoring) és akadémiai integritás}
% Proctoring típusai (élő, rögzített, automatizált), böngészőalapú megoldások,
% adatvédelmi és etikai kérdések. Források: \cite{atoum2017automated,nigam2021systematic,coghlan2021proctor}
\todo{Szöveg}

\section{Meglévő rendszerek összehasonlítása}
% Táblázatos összevetés + szöveges értékelés (saját vélemény!). A táblázatra hivatkozni kell:
A meglévő rendszerek főbb jellemzőit a \ref{tab:osszehasonlitas}.~táblázat foglalja össze.

\begin{table}[H]
  \centering
  \caption{Online vizsgáztató rendszerek összehasonlítása \todo{ellenőrizendő}}
  \label{tab:osszehasonlitas}
  \begin{tabularx}{\textwidth}{@{}l>{\centering\arraybackslash}X>{\centering\arraybackslash}X>{\centering\arraybackslash}X>{\centering\arraybackslash}X@{}}
    \toprule
    Szempont                  & Moodle Quiz & CodeRunner & \todo{?} & Saját rendszer \\
    \midrule
    Vegyes feladattípusok     & igen        & részben    &          & igen \\
    Kódfuttatás sandboxban    & nem         & igen       &          & igen (Docker) \\
    Rejtett tesztesetek       & --          & igen       &          & igen \\
    Valós idejű felügyelet    & nem         & nem        &          & igen \\
    \bottomrule
  \end{tabularx}
\end{table}
